% ================================================================
% 1. INTRODUZIONE
% ================================================================

\section{Introduzione}

\subsection{Scopo del Documento}

Il presente documento definisce le norme, le convenzioni e le procedure operative adottate dal gruppo \textit{Synergo Unit} nello svolgimento del progetto didattico del corso di Ingegneria del Software.

Esso costituisce il riferimento operativo interno del gruppo e ha lo scopo di garantire uniformità, tracciabilità e qualità durante l'intero ciclo di vita del progetto.

Le norme qui descritte sono organizzate in conformità allo standard \textbf{ISO/IEC 12207:1995\textsubscript{G}}, applicato secondo il processo di tailoring\textsubscript{G} descritto nella Sezione~\ref{sec:tailoring}.

Le Norme di Progetto\textsubscript{G} hanno i seguenti obiettivi:
\begin{itemize}
    \item definire il \textit{way of working\textsubscript{G}} del gruppo;
    \item standardizzare le attività svolte;
    \item garantire tracciabilità e verificabilità degli artefatti prodotti;
    \item assicurare uniformità documentale e organizzativa;
    \item supportare il miglioramento continuo del processo;
    \item favorire una collaborazione strutturata tra i membri del gruppo.
\end{itemize}

Il documento è soggetto ad aggiornamento incrementale secondo il principio \textit{just-in-time\textsubscript{G}}: ogni nuova attività deve essere normata prima della sua adozione operativa e ogni evoluzione significativa del progetto può comportare l'introduzione di nuovi processi o l'estensione di quelli esistenti.

Le presenti norme prevalgono su eventuali convenzioni implicite o pratiche informali adottate dal gruppo.

Tutti i membri del gruppo sono tenuti a leggere e rispettare il presente documento.

% ------------------------------------------------

\subsection{Scopo del Prodotto}

Il prodotto sviluppato consiste in una piattaforma documentale intelligente basata su tecnologie AI e Large Language Models (LLM\textsubscript{G}), finalizzata al supporto della scrittura e della gestione di note Markdown\textsubscript{G}.

La piattaforma dovrà supportare funzionalità di:
\begin{itemize}
    \item gestione documentale Markdown;
    \item supporto AI alla scrittura;
    \item generazione di suggerimenti contestuali;
    \item interazione con Large Language Models;
    \item gestione strutturata delle note;
    \item interfaccia utente\textsubscript{G} moderna, accessibile e responsive.
\end{itemize}

% ------------------------------------------------

\subsection{Glossario}

I termini tecnici, gli acronimi e le espressioni potenzialmente ambigue utilizzate all'interno della documentazione sono definiti nel documento
\link{https://synergo-unit-unipd.github.io/Second-Brain/docs/RTB/doc_gestionali/doc_interni/glossario/glossario.pdf}
{Glossario}.

La prima occorrenza dei termini presenti nel glossario viene marcata automaticamente tramite pedice\textsubscript{G} mediante uno script Python dedicato descritto nella Sezione~\ref{sec:strumenti-doc}.

% ------------------------------------------------

\subsection{Convenzioni Tipografiche}

All'interno della documentazione vengono adottate le seguenti convenzioni tipografiche:

\begin{itemize}
    \item i termini inglesi e le espressioni tecniche vengono riportati in \textit{corsivo};
    \item i termini definiti nel Glossario sono marcati con il pedice\textsubscript{G} alla prima occorrenza;
    \item i riferimenti a file, branch\textsubscript{G}, comandi e identificativi tecnici vengono riportati tramite font monospace mediante il comando \texttt{\textbackslash texttt\{\}};
    \item i riferimenti interni alla documentazione vengono riportati tramite hyperlink.
\end{itemize}

% ------------------------------------------------

\subsection{Riferimenti}

\subsubsection{Riferimenti Normativi}

\begin{itemize}
    \item Capitolato d'appalto C6:
    \link{https://www.math.unipd.it/~tullio/IS-1/2025/Progetto/C6p.pdf}
    {Zucchetti\textsubscript{G} S.p.A.};

    \item
    \link{https://synergo-unit-unipd.github.io/Second-Brain/docs/ISO_12207-1995.pdf}
    {Standard ISO/IEC 12207:1995} --
    \textit{Information Technology -- Software Life Cycle Processes};

    \item
    \link{https://www.math.unipd.it/~tullio/IS-1/2025/Dispense/PD1.pdf}
    {Regolamento del Progetto Didattico},
    Prof.\ Tullio Vardanega,
    A.A.\ 2025/2026.
\end{itemize}

% ------------------------------------------------

\subsubsection{Riferimenti Informativi}

\begin{itemize}

    \item Materiale didattico del corso di Ingegneria del Software;

    \item
    Documentazione ufficiale \LaTeX{}:
    \url{https://www.latex-project.org/};

    \item
    Documentazione GitHub:
    \url{https://docs.github.com/};

    \item
    \link{https://synergo-unit-unipd.github.io/Second-Brain/docs/RTB/doc_gestionali/doc_esterni/piano_di_progetto/piano_di_progetto.pdf}
    {Piano di Progetto\textsubscript{G}};

    \item
    \link{https://synergo-unit-unipd.github.io/Second-Brain/docs/RTB/doc_gestionali/doc_esterni/piano_di_qualifica/piano_di_qualifica.pdf}
    {Piano di Qualifica\textsubscript{G}};

    \item
    \link{https://synergo-unit-unipd.github.io/Second-Brain/docs/RTB/doc_tecnici/doc_esterni/analisi_dei_requisiti/analisi_dei_requisiti.pdf}
    {Analisi dei Requisiti\textsubscript{G}};

    \item
    \link{https://synergo-unit-unipd.github.io/Second-Brain/docs/RTB/doc_gestionali/doc_interni/analisi_stato_arte/analisi_stato_arte.pdf}
    {Analisi dello Stato dell'Arte};

    \item Verbali interni ed esterni.
\end{itemize}

% ------------------------------------------------

\subsection{Organizzazione del Documento}

Il presente documento è organizzato nelle seguenti sezioni:

\begin{itemize}
      \item \textbf{Tailoring}: descrive l'adattamento dei processi adottato dal gruppo rispetto al contesto progettuale.

      \item \textbf{Processi Primari}: descrive i processi direttamente coinvolti nello sviluppo del prodotto software;

      \item \textbf{Processi di Supporto}: descrive i processi che supportano lo sviluppo, la qualità, la documentazione e la gestione della configurazione\textsubscript{G};

      \item \textbf{Processi Organizzativi}: descrive le modalità organizzative adottate dal gruppo;
\end{itemize}